% =========================================================================
%  TP2 - UART + ALU : El flujo de la informacion
%  Arquitectura de Computadoras - FCEFyN, UNC
%  Compilar con:  xelatex flujo_tp2.tex  (dos veces, por el indice)
% =========================================================================
\documentclass[11pt,a4paper]{article}

% --- Motor y fuentes (xelatex + fontspec por ruta explicita) -------------
\usepackage{fontspec}
\usepackage{polyglossia}
\setdefaultlanguage{spanish}
\setmainfont{texgyrepagella}[
  Path=/usr/share/texmf/fonts/opentype/public/tex-gyre/,
  Extension=.otf, UprightFont=*-regular, BoldFont=*-bold,
  ItalicFont=*-italic, BoldItalicFont=*-bolditalic]
\setsansfont{texgyreheros}[
  Path=/usr/share/texmf/fonts/opentype/public/tex-gyre/,
  Extension=.otf, UprightFont=*-regular, BoldFont=*-bold,
  ItalicFont=*-italic, BoldItalicFont=*-bolditalic]
\setmonofont{texgyrecursor}[
  Path=/usr/share/texmf/fonts/opentype/public/tex-gyre/,
  Extension=.otf, UprightFont=*-regular, BoldFont=*-bold,
  ItalicFont=*-italic, BoldItalicFont=*-bolditalic, Scale=0.92]

\usepackage{amsmath,amssymb}
\usepackage[margin=2.3cm]{geometry}
\usepackage{tikz}
\usetikzlibrary{arrows.meta,positioning,calc,shapes.geometric,
                decorations.pathreplacing,fit,backgrounds}
\usepackage[most]{tcolorbox}
\usepackage{booktabs}
\usepackage{array}
\usepackage{needspace}
\usepackage{enumitem}
\usepackage{fancyhdr}
\usepackage{hyperref}
\hypersetup{colorlinks=true, linkcolor=azul, urlcolor=azul,
            pdftitle={TP2 UART + ALU - Flujo de la informacion}}

\linespread{1.20}

% Red de seguridad para justificado: si el sistema donde se compila no tiene
% los patrones de silabeo del espanol, esto evita lineas desbordadas.
\tolerance=1500
\emergencystretch=2em
\hbadness=4000

% --- Paleta -------------------------------------------------------------
\definecolor{azul}{HTML}{3573C4}
\definecolor{naranja}{HTML}{D97706}
\definecolor{violeta}{HTML}{8B5CF6}
\definecolor{verde}{HTML}{2E8B57}
\definecolor{gris}{HTML}{4B5563}
\definecolor{fondo}{HTML}{FCFCFB}
\pagecolor{fondo}

% --- Cajas --------------------------------------------------------------
\newtcolorbox{intuicion}[1]{
  enhanced, colback=verde!6, colframe=verde, coltitle=white,
  fonttitle=\bfseries\sffamily\small, title={#1},
  boxrule=0.7pt, arc=2pt, left=5pt, right=5pt, top=4pt, bottom=4pt,
  breakable}

\newtcolorbox{ficha}[1]{
  enhanced, colback=azul!5, colframe=azul, coltitle=white,
  fonttitle=\bfseries\sffamily\small, title={#1},
  boxrule=0.7pt, arc=2pt, left=5pt, right=5pt, top=4pt, bottom=4pt,
  breakable}

\newtcolorbox{ojo}[1]{
  enhanced, colback=naranja!7, colframe=naranja, coltitle=white,
  fonttitle=\bfseries\sffamily\small, title={#1},
  boxrule=0.7pt, arc=2pt, left=5pt, right=5pt, top=4pt, bottom=4pt,
  breakable}

\newtcolorbox{clave}{
  enhanced, colback=violeta!7, colframe=violeta,
  boxrule=0.9pt, arc=2pt, left=6pt, right=6pt, top=5pt, bottom=5pt,
  breakable}

% --- Titulos ------------------------------------------------------------
\usepackage{titlesec}
\titleformat{\section}{\Large\bfseries\sffamily\color{azul}}{\thesection}{0.6em}{}
\titleformat{\subsection}{\large\bfseries\sffamily\color{gris}}{\thesubsection}{0.6em}{}
\titlespacing*{\section}{0pt}{16pt}{8pt}

\pagestyle{fancy}
\fancyhf{}
\fancyhead[L]{\sffamily\footnotesize\color{gris} TP2 --- UART + ALU}
\fancyhead[R]{\sffamily\footnotesize\color{gris} El flujo de la informaci\'on}
\fancyfoot[C]{\sffamily\footnotesize\color{gris}\thepage}
\renewcommand{\headrulewidth}{0.4pt}

\newcommand{\sg}[1]{\texttt{\small #1}}

% =========================================================================
\begin{document}

% ---------------------------- PORTADA -----------------------------------
\begin{titlepage}
\centering
\vspace*{2.5cm}
{\sffamily\Huge\bfseries\color{azul} El flujo de la informaci\'on}\\[0.4cm]
{\sffamily\LARGE\color{gris} TP2 --- UART + ALU}\\[1.2cm]
{\large Del byte que sale de la PC al byte que vuelve,\\
siguiendo el caminito bloque por bloque}\\[2cm]

\begin{tcolorbox}[width=0.82\textwidth, colback=white, colframe=gris,
                  boxrule=0.7pt, arc=3pt]
\sffamily\small
\textbf{\color{azul} C\'omo leer este apunte}\\[4pt]
\textcolor{verde}{$\blacksquare$} \textbf{Verde:} la intuici\'on. La analog\'ia antes de la f\'ormula.\\[2pt]
\textcolor{azul}{$\blacksquare$} \textbf{Azul:} la ficha t\'ecnica. Se\~nales, n\'umeros y ciclos exactos.\\[2pt]
\textcolor{naranja}{$\blacksquare$} \textbf{Naranja:} lo que conviene mirar dos veces, o preguntar en la c\'atedra.\\[2pt]
\textcolor{violeta}{$\blacksquare$} \textbf{Violeta:} la idea central de la secci\'on. Si le\'es solo esto, ya ganaste algo.
\end{tcolorbox}

\vfill
{\sffamily\small\color{gris} Arquitectura de Computadoras --- FCEFyN, UNC \\
Basys 3 (Artix-7) --- clock de 100\,MHz --- 19200 baud}
\end{titlepage}

% ------------------- EL TP EN UNA PAGINA --------------------------------
\section*{El TP2 entero en una p\'agina}
\addcontentsline{toc}{section}{El TP2 entero en una p\'agina}

La PC manda tres bytes por un cable: el dato \textbf{A}, el dato \textbf{B} y el
\textbf{c\'odigo de operaci\'on}. La placa contesta un byte: el \textbf{resultado}.
Eso es todo lo que se ve desde afuera.

Adentro pasa lo siguiente. El cable trae \emph{un bit por vez} y no trae reloj:
llega un \emph{nivel} que sube y baja, y la FPGA tiene que adivinar d\'onde
empieza y termina cada bit. Para adivinar bien usa un truco: cuenta mucho m\'as
r\'apido que el cable. Genera un pulso llamado \sg{s\_tick} \textbf{16 veces por
cada bit}, y con esos 16 pulsos se ubica justo en el medio del bit, donde la
se\~nal est\'a m\'as estable, y ah\'i la mira.

Una vez que junt\'o los 8 bits, el byte entra en una cola (\sg{fifo\_rx}) y sale
del mundo lento del cable para entrar al mundo r\'apido de los 100\,MHz. Ah\'i lo
levanta una m\'aquina de estados (\sg{intf}) que va llenando tres registros:
primero A, despu\'es B, despu\'es OP. Cuando tiene los tres, la ALU --- que es
combinacional pura, sin reloj --- ya tiene el resultado en su salida. La FSM lo
mete en otra cola (\sg{fifo\_tx}), el transmisor lo saca bit por bit al mismo
ritmo de 16 ticks por bit, y el byte vuelve a la PC.

\begin{clave}
\textbf{La idea que hay que llevarse:} en este sistema conviven
\textbf{tres velocidades} muy distintas --- el clock (10\,ns), el tick
(3{,}26\,\textmu s) y el bit (52{,}16\,\textmu s) --- y \emph{cada bloque vive en
una de ellas}. Las colas FIFO existen justamente para que un bloque r\'apido y
uno lento puedan hablarse sin perder datos. Entender qui\'en vive en qu\'e
velocidad es entender el TP.
\end{clave}

\newpage
\tableofcontents
\newpage

% =========================================================================
\section{El vocabulario m\'inimo}

Seis palabras. Si estas seis est\'an claras, el resto del apunte se lee solo.

\begin{ficha}{Las seis palabras}
\begin{description}[leftmargin=0pt, labelindent=0pt, labelsep=0.45em,
                    itemsep=5pt, parsep=0pt, topsep=2pt,
                    font=\bfseries\sffamily\color{azul}]
\item[Baud rate.] Bits por segundo que viajan por el cable. Ac\'a: 19200. Es la
  velocidad del \emph{cable}, no la del chip.
\item[Trama (frame).] El paquetito completo que viaja por el cable: 1 bit de
  start (siempre 0) $+$ 8 bits de datos $+$ 1 bit de stop (siempre 1). 10 bits
  en total para transportar 8.
\item[Sobremuestreo.] Mirar la l\'inea muchas veces por bit. Ac\'a, 16. Sirve
  para ubicar el centro del bit con precisi\'on de $1/16$ de bit.
\item[Tick (\texttt{s\_tick}).] Un pulso de \textbf{un solo ciclo de clock} de
  ancho, que aparece 16 veces por bit. Es el metr\'onomo del UART.
\item[Pulso de un ciclo.] Una se\~nal que vale 1 durante exactamente un ciclo de
  clock (\sg{rx\_done\_tick}, \sg{tx\_done\_tick}). Sirve para avisar
  \emph{``pas\'o algo''} sin trabar nada.
\item[FIFO.] Una cola. Se escribe por un lado, se lee por el otro, y los dos
  lados pueden ir a velocidades distintas.
\end{description}
\end{ficha}

\begin{intuicion}{Por qu\'e ``as\'incrono''}
El cable tiene \textbf{una sola l\'inea de datos y ning\'un reloj}. Es como
alguien que te habla desde la otra punta de la cancha sin que ustedes dos
hayan acordado un metr\'onomo: vos sab\'es \emph{a qu\'e velocidad} habla
(los 19200 baud, pactados de antemano), pero no sab\'es \emph{cu\'ando} va a
empezar la pr\'oxima palabra. Por eso cada trama arranca con un bit de start:
es el ``ey, escuchame'' que dispara tu cronometraje.
\end{intuicion}

% =========================================================================
\section{El mapa: qui\'en le habla a qui\'en}

Antes de recorrer el camino, la foto completa. Las flechas gruesas son datos
(8 bits), las finas son se\~nales de control (1 bit).

\vspace{4pt}
\begin{center}
\begin{tikzpicture}[
  scale=0.97, transform shape, font=\sffamily\scriptsize,
  blk/.style={draw=azul!70, rounded corners=2.5pt, minimum height=1.05cm,
              minimum width=2.0cm, align=center, fill=azul!8, line width=0.8pt},
  ff/.style={blk, draw=verde!70, fill=verde!10},
  bg/.style={blk, draw=naranja!80, fill=naranja!12},
  dato/.style={-{Stealth[length=2.4mm]}, line width=1.6pt, draw=gris!85},
  ctrl/.style={-{Stealth[length=2mm]}, line width=0.7pt, draw=naranja, densely dashed},
  lbl/.style={font=\sffamily\tiny, inner sep=1.5pt},
  sig/.style={font=\sffamily\tiny, text=gris, inner sep=1.5pt}
]

% --- bloques ---
\node[blk] (rx)   at (0,2.3)   {\texttt{uart\_rx}};
\node[ff]  (frx)  at (3.4,2.3) {\texttt{fifo\_rx}};
\node[blk, minimum height=1.6cm] (intf) at (7.0,1.15) {\texttt{intf}\\[-1pt]\tiny FSM};
\node[blk, minimum height=1.6cm] (alu)  at (10.3,1.15) {\texttt{ALU}\\[-1pt]\tiny combinacional};
\node[ff]  (ftx)  at (3.4,0)   {\texttt{fifo\_tx}};
\node[blk] (tx)   at (0,0)     {\texttt{uart\_tx}};
\node[bg, minimum height=1.1cm] (bgen) at (-3.3,1.15) {\texttt{baud\_gen}\\[-1pt]\tiny m\'odulo 326};

% --- linea serie ---
\draw[dato] (-1.55,2.3) -- (rx.west);
\draw[dato] (tx.west) -- (-1.55,0);
\node[sig, above] at (-1.3,2.38) {de la PC};
\node[sig, below] at (-1.3,-0.08) {a la PC};

% --- s_tick ---
\draw[ctrl] (bgen.east) -- (-2.0,1.15) -- (-2.0,1.72) -- (rx.south);
\draw[ctrl] (bgen.east) -- (-2.0,1.15) -- (-2.0,0.58) -- (tx.north);
\node[sig, text=naranja, anchor=west] at (-1.95,1.15) {\texttt{s\_tick}};

% --- rx -> fifo_rx ---
\draw[dato] (rx.east) -- (frx.west);
\node[sig, above] at (1.7,2.42) {\texttt{dout}};
\node[sig, below] at (1.7,2.18) {\texttt{rx\_done\_tick}};

% --- fifo_rx -> intf (datos) ---
\draw[dato] (frx.east) -- (5.15,2.3) -- (5.15,1.55) -- (intf.west |- 0,1.55);
\node[sig, above] at (4.72,2.38) {\texttt{r\_data}};
\node[sig, above] at (5.72,1.58) {\texttt{rx\_empty}};
% --- intf -> fifo_rx (rd) ---
\draw[ctrl] (intf.west |- 0,0.75) -- (5.62,0.75) -- (5.62,2.02) -- (frx.east |- 0,2.02);
\node[sig, text=naranja, below] at (4.72,2.0) {\texttt{rd}};

% --- intf <-> alu ---
\draw[dato] (intf.east |- 0,1.5) -- (alu.west |- 0,1.5);
\node[sig, above] at (8.65,1.53) {\texttt{a, b, op}};
\draw[dato] (alu.west |- 0,0.8) -- (intf.east |- 0,0.8);
\node[sig, below] at (8.65,0.77) {\texttt{result, carry}};

% --- intf -> fifo_tx ---
\draw[dato] (intf.south) -- (7.0,-1.05) -- (3.4,-1.05) -- (ftx.south);
\node[sig, below] at (5.2,-1.05) {\texttt{w\_data} \ / \ \texttt{wr} \ / \ \texttt{tx\_full}};

% --- fifo_tx -> tx ---
\draw[dato] (ftx.west) -- (tx.east);
\node[sig, above] at (1.7,0.12) {\texttt{din} / \texttt{tx\_start}};
\node[sig, below] at (1.7,-0.12) {\texttt{tx\_done\_tick}};

% --- leyenda de mundos ---
\draw[draw=naranja!45, densely dotted, line width=0.8pt, rounded corners]
  (-1.35,-0.72) rectangle (1.3,3.05);
\node[sig, text=naranja, align=center] at (-0.02,3.28) {mundo lento\\[-2pt]ritmo del \texttt{s\_tick}};
\draw[draw=azul!45, densely dotted, line width=0.8pt, rounded corners]
  (5.75,-0.72) rectangle (11.45,2.15);
\node[sig, text=azul, align=center] at (8.6,2.38) {mundo r\'apido --- ritmo del \texttt{clk}};
\node[sig, text=verde!55!black, align=center] at (3.4,3.28) {frontera\\[-2pt]de velocidades};

\end{tikzpicture}
\end{center}
\vspace{-2pt}

\begin{ficha}{El reparto de tareas}
\renewcommand{\arraystretch}{1.15}
\begin{tabular}{@{}>{\ttfamily\small}l p{5.1cm} p{4.5cm}@{}}
\toprule
\textnormal{\textbf{\sffamily Bloque}} & \textbf{\sffamily Qu\'e hace} & \textbf{\sffamily A qu\'e ritmo vive} \\
\midrule
baud\_gen & Cuenta 326 ciclos y tira un tick & clock (100\,MHz) \\
uart\_rx  & Cuenta ticks y arma el byte & tick (306\,kHz) \\
fifo\_rx  & Frontera lento $\to$ r\'apido & los dos lados \\
intf      & Ordena A, B, OP y pide el resultado & clock (100\,MHz) \\
ALU       & Calcula. No tiene reloj. & instant\'aneo (comb.) \\
fifo\_tx  & Frontera r\'apido $\to$ lento & los dos lados \\
uart\_tx  & Cuenta ticks y saca el byte & tick (306\,kHz) \\
\bottomrule
\end{tabular}
\end{ficha}

% =========================================================================
\section{Las tres escalas de tiempo}
\label{sec:escalas}

Esta es la secci\'on que hay que entender antes que ninguna otra. Todo el resto
del apunte se apoya ac\'a.

\begin{intuicion}{La analog\'ia del tren}
Pens\'a en tres relojes encajados uno adentro del otro:
\begin{itemize}[leftmargin=1.2em, itemsep=1pt]
\item El \textbf{clock} es el segundero de un cron\'ometro: tic cada 10\,ns.
  Todo lo que es \sg{always @(posedge clk)} se mueve con \'el.
\item El \textbf{tick} es una campanita que suena \textbf{cada 326 tics} del
  cron\'ometro. El \sg{baud\_gen} es el \'unico que la hace sonar.
\item El \textbf{bit} en el cable dura \textbf{16 campanitas}.
\end{itemize}
El clock es el \'unico reloj f\'isico. El tick y el bit son ficciones que
construimos contando: el tick contando ciclos de clock, el bit contando ticks.
\end{intuicion}

\vspace{2pt}
\begin{center}
\begin{tikzpicture}[font=\sffamily\scriptsize, x=1cm, y=1cm]
% nivel bit
\draw[fill=azul!12, draw=azul!60, line width=0.6pt] (0,2.1) rectangle (13,2.75);
\node at (6.5,2.42) {\textbf{1 bit en el cable} = 52,16 \textmu s = 5216 ciclos de clock};
% nivel tick
\foreach \i in {0,...,15} {
  \draw[fill=verde!12, draw=verde!60, line width=0.5pt]
    (\i*0.8125,1.25) rectangle (\i*0.8125+0.78,1.85);
}
\node[align=center] at (6.5,0.95) {\textcolor{verde}{16 ticks}, uno cada 3,26 \textmu s (326 ciclos de clock)};
% nivel clock
\draw[fill=naranja!12, draw=naranja!70, line width=0.5pt] (0,0.15) rectangle (0.78,0.6);
\node[font=\sffamily\tiny] at (0.39,0.375) {326};
\node[anchor=west, align=left] at (1.0,0.375)
  {\textcolor{naranja}{326 ciclos de clock} de 10 ns cada uno $=$ 1 tick};
% llaves
\draw[decorate, decoration={brace, amplitude=3pt}, draw=verde!70]
  (0,1.2) -- (0.78,1.2);
\draw[->, draw=verde!70, line width=0.5pt] (0.39,1.15) -- (0.39,0.68);
\end{tikzpicture}
\end{center}

\begin{ficha}{Los n\'umeros, de una vez}
\[
M \;=\; \frac{f_{clk}}{\text{BaudRate} \times 16}
   \;=\; \frac{100\,000\,000}{19200 \times 16}
   \;=\; 325{,}52 \;\approx\; \mathbf{326}
\]
\renewcommand{\arraystretch}{1.2}
\begin{tabular}{@{}l r r l@{}}
\toprule
\textbf{\sffamily Escala} & \textbf{\sffamily Per\'iodo} & \textbf{\sffamily Frecuencia} & \textbf{\sffamily Qui\'en la genera} \\
\midrule
Ciclo de clock        & 10 ns          & 100 MHz    & el cristal de la placa \\
Tick (\sg{s\_tick})   & 3,26 \textmu s & 306,7 kHz  & \sg{baud\_gen} contando 326 \\
Bit en el cable       & 52,16 \textmu s& 19,18 kbaud& \sg{uart\_rx}/\sg{uart\_tx} contando 16 ticks \\
Trama de 10 bits      & 521,6 \textmu s& ---        & la trama completa \\
\bottomrule
\end{tabular}
\end{ficha}

\begin{ojo}{El error de redondeo, y por qu\'e no molesta}
326 no es 325,52. El bit real dura 52,16 \textmu s en vez de 52,083:
un error de \textbf{+0,147 por ciento}.

\smallskip
$\rightarrow$ \textbf{¿Alcanza?} El receptor se resincroniza en cada bit de start,
as\'i que el error solo se acumula durante \textbf{una trama}. El \'ultimo bit que
mira (el de stop) lo muestrea 9,5 bits despu\'es del flanco de start. Para no
equivocarse, el corrimiento acumulado tiene que ser menor a medio bit:
\[
\text{tolerancia} \;=\; \frac{0{,}5}{9{,}5} \;=\; 5{,}26 \text{ por ciento}
\]
Tenemos 0,147. Estamos 35 veces por debajo del l\'imite. Por eso no hace falta
ninguna correcci\'on fraccionaria.
\end{ojo}

% =========================================================================
\section{El viaje de ida: de la PC a la ALU}

Ahora s\'i, el caminito. Vamos a seguir \textbf{un solo byte} --- digamos el
dato A $=$ \texttt{0x05} --- desde que sale de la PC hasta que queda guardado en
un registro adentro de la FPGA.

\subsection{Paso 0 --- La PC suelta el byte}

La PC pone en el cable, en orden y a 19200 bits por segundo:

\vspace{4pt}
\begin{center}
\begin{tikzpicture}[font=\sffamily\scriptsize, x=1.05cm, y=0.62cm]
\foreach \x/\t/\c in {0/start/naranja, 1/1/azul, 2/0/azul, 3/1/azul, 4/0/azul,
                      5/0/azul, 6/0/azul, 7/0/azul, 8/0/azul, 9/stop/verde} {
  \draw[fill=\c!12, draw=\c!70, line width=0.6pt] (\x,0) rectangle (\x+0.96,0.85);
  \node at (\x+0.48,0.42) {\t};
}
\foreach \x/\b in {1/D0,2/D1,3/D2,4/D3,5/D4,6/D5,7/D6,8/D7} {
  \node[font=\sffamily\tiny, gris] at (\x+0.48,-0.35) {\b};
}
\node[font=\sffamily\tiny] at (0.48,-0.35) {0};
\node[font=\sffamily\tiny] at (9.48,-0.35) {1};
\draw[decorate, decoration={brace, amplitude=4pt}, draw=gris]
  (9.96,1.05) -- (0,1.05);
\node[above, font=\sffamily\tiny] at (5,1.2) {10 bits $=$ 521,6 \textmu s};
\end{tikzpicture}
\end{center}

\begin{ojo}{El bit menos significativo va primero}
\texttt{0x05} $=$ \texttt{0000\,0101} en binario, pero por el cable sale
\texttt{1010\,0000}: \textbf{D0 primero}. Es por esto que en \sg{uart\_rx} el
shift register se arma metiendo el bit nuevo \emph{por izquierda}:
\begin{center}\sg{b\_next = \{rx, b\_reg[DBIT-1:1]\};}\end{center}
Cada bit nuevo empuja a los anteriores hacia la derecha. Cuando entr\'o el octavo,
el primero que lleg\'o qued\'o en el bit 0. Si esto se hace al rev\'es el byte sale
espejado y es uno de los bugs cl\'asicos del TP.
\end{ojo}

\subsection{Paso 1 --- \texttt{baud\_gen}: el metr\'onomo que nunca para}

El \sg{baud\_gen} no espera nada ni le importa lo que pase en el cable. Arranca
con el reset y cuenta para siempre: $0, 1, 2, \dots, 325, 0, 1, \dots$ Cuando el
contador vale 325 pone \sg{tick} en 1 durante \textbf{un solo ciclo de clock}, y
en el siguiente vuelve a 0.

\begin{ficha}{La l\'inea que importa}
\begin{center}\sg{assign tick = (count\_reg == (M - 1));}\end{center}
Es combinacional sobre el valor \emph{actual} del contador, no un registro. Por eso
el pulso dura exactamente un ciclo: al ciclo siguiente el contador ya vale 0 y la
comparaci\'on da falso. Esto es lo que hace que \sg{s\_tick} sirva como
\emph{habilitaci\'on} (\emph{clock enable}) y no como un segundo reloj.
\end{ficha}

\begin{intuicion}{Por qu\'e no se usa el tick como clock}
Tentaci\'on t\'ipica: ``si necesito algo que vaya a 306\,kHz, ¿por qu\'e no hago
\sg{always @(posedge tick)}?''. Porque eso crea un \textbf{segundo dominio de
reloj} dentro del chip, con todos los problemas de cruce de dominios que eso
trae, y adem\'as en una FPGA una se\~nal generada por l\'ogica no tiene la red de
distribuci\'on de reloj: llega con \emph{skew} a cada flip-flop.

La forma correcta es la que usa el TP: \textbf{todo corre con el mismo clock de
100\,MHz}, y el tick se usa como una condici\'on (\sg{if (s\_tick) ...}) que
habilita el avance. As\'i hay un solo dominio de reloj en todo el dise\~no.
\end{intuicion}

\subsection{Paso 2 --- El sincronizador: dos flip-flops antes de tocar nada}

La l\'inea \sg{i\_rx} viene de un chip externo y cambia cuando se le canta. No
tiene ninguna relaci\'on con el clock de 100\,MHz. Si la met\'es directo a un
flip-flop y justo cambia en el instante del flanco, el flip-flop puede quedar
\textbf{metaestable}: ni 0 ni 1, oscilando un rato antes de decidirse.

\begin{ficha}{Los dos flip-flops del top}
\begin{center}
\sg{rx\_meta <= i\_rx;} \qquad \sg{rx\_sync <= rx\_meta;}
\end{center}
El primero es el que se banca la metaestabilidad; el segundo le da un ciclo entero
para que se estabilice antes de que el resto del dise\~no lo vea. Costo:
\textbf{20\,ns de retardo}. Contra un bit de 52\,160\,ns, es 0,04 por ciento.
Gratis.

\smallskip
En el \sg{.xdc} esto se acompa\~na con \sg{set\_false\_path -from [get\_ports i\_rx]}:
le avisa a Vivado que no analice el timing de esa entrada, porque ya sabemos que
es as\'incrona y ya la tratamos.
\end{ficha}

\subsection{Paso 3 --- \texttt{uart\_rx}: contar ticks para encontrar el centro}

Ac\'a est\'a el coraz\'on del TP. La FSM tiene cuatro estados y dos contadores:

\begin{itemize}[leftmargin=1.3em, itemsep=2pt]
\item \sg{s\_reg}: cuenta \textbf{ticks} dentro del bit actual (0 a 15).
\item \sg{n\_reg}: cuenta \textbf{bits de datos} ya recibidos (0 a 7).
\end{itemize}

\vspace{4pt}
\begin{center}
\begin{tikzpicture}[
  font=\sffamily\scriptsize,
  st/.style={circle, draw=azul!70, fill=azul!8, minimum size=1.35cm,
             line width=0.8pt, align=center},
  ar/.style={-{Stealth[length=2mm]}, line width=0.7pt, draw=gris},
  lbl/.style={font=\sffamily\tiny, fill=fondo, inner sep=1.5pt}
]
\node[st] (idle)  at (0,0)    {IDLE};
\node[st] (start) at (3.5,0)  {START};
\node[st] (data)  at (7,0)    {DATA};
\node[st] (stop)  at (10.5,0) {STOP};

\draw[ar] (idle) -- (start) node[lbl, midway, above] {\texttt{rx==0}}
                            node[lbl, midway, below] {$s \leftarrow 0$};
\draw[ar] (start) -- (data)  node[lbl, midway, above] {\texttt{s==7}}
                            node[lbl, midway, below] {$s,n \leftarrow 0$};
\draw[ar] (data) -- (stop)   node[lbl, midway, above] {\texttt{s==15 \&\& n==7}};
\draw[ar] (stop.north) .. controls +(0,1.5) and +(0,1.5) .. (idle.north)
  node[lbl, midway, above] {\texttt{s==SB\_TICK-1} $\Rightarrow$ \texttt{rx\_done\_tick=1}};

\path[ar] (idle)  edge[loop below, looseness=7, min distance=8mm]
  node[lbl, below=2pt] {\texttt{rx==1}} (idle);
\path[ar] (start) edge[loop below, looseness=7, min distance=8mm]
  node[lbl, below=2pt] {\texttt{s\_tick} $\Rightarrow s \leftarrow s+1$} (start);
\path[ar] (data)  edge[loop below, looseness=7, min distance=8mm]
  node[lbl, below=2pt] {\texttt{s\_tick} $\Rightarrow s{+}1$, y cada 16 \, $n{+}1$} (data);
\path[ar] (stop)  edge[loop below, looseness=7, min distance=8mm]
  node[lbl, below=2pt] {\texttt{s\_tick} $\Rightarrow s \leftarrow s+1$} (stop);
\end{tikzpicture}
\end{center}

\begin{clave}
\textbf{Los dos n\'umeros m\'agicos son 7 y 15.} En \sg{START} se esperan
\textbf{8 ticks} (contar hasta 7 desde 0) porque 8 es \emph{medio bit}: eso deja
al receptor parado en el centro del bit de start. A partir de ah\'i, en
\sg{DATA} se esperan \textbf{16 ticks} (contar hasta 15) por cada bit, que es
\emph{un bit entero}. Resultado: cada muestreo cae en el centro del bit
correspondiente.
\end{clave}

\subsubsection*{D\'onde cae cada muestra}

Esto es lo que est\'a pasando en el cable mientras la FSM cuenta:

\vspace{4pt}
\begin{center}
\begin{tikzpicture}[x=0.245cm, y=1cm, font=\sffamily\tiny]
\def\hi{0.78}
% señal rx: idle=1, start=0 (0..16), D0=1 (16..32), D1=0 (32..48), D2=1(48..56)
\draw[line width=1.1pt, draw=azul]
  (-6,\hi) -- (0,\hi) -- (0,0) -- (16,0) -- (16,\hi) -- (32,\hi)
  -- (32,0) -- (48,0) -- (48,\hi) -- (56,\hi);
% marcas de tick
\foreach \t in {0,...,56} {
  \draw[draw=verde!55, line width=0.35pt] (\t,-0.30) -- (\t,-0.14);
}
\foreach \t in {0,16,32,48} {
  \draw[draw=gris!45, densely dotted, line width=0.5pt] (\t,-0.32) -- (\t,1.05);
}
% etiquetas de bit
\node[above] at (-3,\hi) {reposo};
\node[above] at (8,\hi+0.02) {\textbf{start}};
\node[above] at (24,\hi+0.02) {\textbf{D0}};
\node[above] at (40,\hi+0.02) {\textbf{D1}};
% flechas de muestreo
\foreach \x/\lb in {8/{s=7}, 24/{s=15}, 40/{s=15}} {
  \draw[-{Stealth[length=1.8mm]}, draw=naranja, line width=0.9pt]
    (\x,-0.72) -- (\x,-0.06);
  \node[below, naranja] at (\x,-0.70) {\lb};
}
\node[below, naranja, align=center] at (8,-1.0) {centro del\\bit de start};
\node[below, naranja, align=center] at (24,-1.0) {muestrea\\y guarda};
\node[below, naranja, align=center] at (40,-1.0) {muestrea\\y guarda};
% llave 16 ticks (arriba, para no pisar las etiquetas de muestreo)
\draw[decorate, decoration={brace, amplitude=3.5pt}, draw=verde!70, line width=0.6pt]
  (16,1.28) -- (32,1.28);
\node[above, verde!60!black] at (24,1.32) {16 ticks $=$ 1 bit $=$ 52,16 \textmu s};
% tick label
\node[verde!60!black, anchor=west] at (57,-0.22) {\texttt{s\_tick}};
\end{tikzpicture}
\end{center}
\vspace{6pt}

\begin{ficha}{La cuenta exacta, en ticks desde el flanco de start}
\renewcommand{\arraystretch}{1.15}
\begin{tabular}{@{}l r r l@{}}
\toprule
\textbf{\sffamily Qu\'e} & \textbf{\sffamily Tick n\textsuperscript{o}} & \textbf{\sffamily En bits} & \textbf{\sffamily Estado} \\
\midrule
Detecta \sg{rx==0}      & 0   & 0,0 & \sg{IDLE} $\to$ \sg{START} \\
Centro del bit de start & 8   & 0,5 & \sg{START} $\to$ \sg{DATA} \\
Muestrea D0             & 24  & 1,5 & \sg{DATA} \\
Muestrea D1             & 40  & 2,5 & \sg{DATA} \\
\multicolumn{4}{@{}l}{\quad\dots\ un bit cada 16 ticks \dots} \\
Muestrea D7             & 136 & 8,5 & \sg{DATA} $\to$ \sg{STOP} \\
Fin del bit de stop     & 152 & 9,5 & \sg{STOP} $\to$ \sg{IDLE}, \sg{rx\_done\_tick=1} \\
\bottomrule
\end{tabular}
\end{ficha}

\begin{ojo}{Detalle fino: el flanco de start no est\'a alineado con el tick}
El \sg{baud\_gen} corre libre, as\'i que cuando llega el flanco de bajada de
\sg{rx} el contador puede estar en cualquier valor entre 0 y 325. Es decir: el
primer tick despu\'es de entrar en \sg{START} llega en un momento aleatorio,
entre 0 y 3,26 \textmu s despu\'es.

\smallskip
Eso mete un error de hasta \textbf{un tick $=$ 1/16 de bit $=$ 6,25 por ciento}
en la ubicaci\'on del punto de muestreo. Sumado al 0,147 por ciento de deriva
acumulada, el peor caso queda en $\sim$8 por ciento de bit de corrimiento
respecto del centro ideal --- todav\'ia muy lejos del 50 por ciento que ser\'ia
un error. \textbf{Este es el motivo real por el que se sobremuestrea 16 veces y
no 4:} con 4 ticks por bit, la incertidumbre ser\'ia de 25 por ciento y el margen
se pondr\'ia inc\'omodo.
\end{ojo}

\subsection{Paso 4 --- \texttt{rx\_done\_tick}: el aviso de un ciclo}

Cuando la FSM termina el bit de stop levanta \sg{rx\_done\_tick} durante
\textbf{un ciclo de clock}. No es un flag que queda prendido: es un pulso.

\begin{intuicion}{Por qu\'e un pulso y no un flag}
Un flag habr\'ia que apagarlo, y apagarlo requiere que alguien avise que ya lo
ley\'o --- lo que trae un ida y vuelta de handshake. Un pulso de un ciclo se
autolimpia: quien lo tiene que ver lo ve, porque todos miran el mismo clock.
Ac\'a el pulso va directo a la entrada \sg{wr} de la FIFO, as\'i que el efecto es
``escrib\'i este byte, ahora''.
\end{intuicion}

\subsection{Paso 5 --- \texttt{fifo\_rx}: la frontera entre los dos mundos}

Este es el punto exacto donde el byte cambia de velocidad.

\begin{ficha}{Qu\'e pasa, ciclo por ciclo}
\renewcommand{\arraystretch}{1.2}
\begin{tabular}{@{}>{\bfseries\sffamily}l p{11.2cm}@{}}
\toprule
Ciclo $N$   & \sg{uart\_rx} est\'a en \sg{STOP} con \sg{s\_reg==15} y llega
              \sg{s\_tick}. Combinacionalmente: \sg{rx\_done\_tick = 1}. \\
Flanco $N{\to}N{+}1$ & La FIFO escribe el byte en \sg{array\_reg[w\_ptr]},
              avanza \sg{w\_ptr} y baja \sg{empty\_reg}. \\
Ciclo $N{+}1$ & \sg{rx\_empty == 0}. El \sg{intf} lo ve en el mismo ciclo, porque
              es combinacional sobre \sg{empty\_reg}. Y \sg{r\_data} ya muestra el
              byte, sin esperar. \\
\bottomrule
\end{tabular}

\smallskip
La FIFO es \emph{first-word fall-through}: \sg{assign r\_data = array\_reg[r\_ptr\_reg];}
es una lectura combinacional. No hay que pedir el dato y esperar un ciclo: el dato
de la cabeza de la cola \textbf{siempre} est\'a en la salida.
\end{ficha}

\begin{ojo}{El truco de los dos flags}
\sg{w\_ptr} y \sg{r\_ptr} tienen el mismo ancho (4 bits), as\'i que cuando son
iguales la cola puede estar \textbf{vac\'ia o llena} --- son el mismo estado de los
punteros. Por eso el m\'odulo lleva dos registros aparte, \sg{full\_reg} y
\sg{empty\_reg}, en vez de deducirlo comparando punteros. El \sg{case (\{wr, rd\})}
es lo que los mantiene: escribir baja \sg{empty}, leer baja \sg{full}, y solo si
el puntero que avanza alcanza al otro se levanta el flag correspondiente.
\end{ojo}

\begin{intuicion}{¿Y para qu\'e sirve la cola, si el \texttt{intf} es much\'isimo
m\'as r\'apido?}
Dos razones concretas.

\smallskip
\textbf{Una:} el \sg{intf} podr\'ia estar ocupado en \sg{SEND} esperando a que
se libere la FIFO de transmisi\'on justo cuando llega un byte nuevo. Sin cola,
ese byte se perder\'ia (eso se llama \emph{overrun}).

\smallskip
\textbf{Dos:} la cola es lo que permite que la PC dispare las tres tramas
seguidas sin esperar respuesta. Llegan tres bytes en 1,56\,ms y quedan
guardados; el \sg{intf} los va sacando a su ritmo. Con 16 lugares de
profundidad hay margen de sobra.
\end{intuicion}

\subsection{Paso 6 --- \texttt{intf}: armar la operaci\'on}

La FSM del \sg{intf} vive en el mundo r\'apido: un estado por ciclo de clock.
Su trabajo es trivial en comparaci\'on con lo que tard\'o el byte en llegar.

\vspace{4pt}
\begin{center}
\begin{tikzpicture}[
  font=\sffamily\scriptsize,
  st/.style={rounded corners=3pt, draw=azul!70, fill=azul!8, minimum width=2.3cm,
             minimum height=1.15cm, line width=0.8pt, align=center},
  ar/.style={-{Stealth[length=2mm]}, line width=0.7pt, draw=gris},
  lbl/.style={font=\sffamily\tiny, fill=fondo, inner sep=1.5pt}
]
\node[st] (a)  at (0,0)    {LOAD\_A\\[-1pt]\tiny $a \leftarrow$ \texttt{r\_data}};
\node[st] (b)  at (3.7,0)  {LOAD\_B\\[-1pt]\tiny $b \leftarrow$ \texttt{r\_data}};
\node[st] (op) at (7.4,0)  {LOAD\_OP\\[-1pt]\tiny $op \leftarrow$ \texttt{r\_data}};
\node[st, fill=verde!10, draw=verde!70] (s) at (11.1,0)
  {SEND\\[-1pt]\tiny \texttt{w\_data = result}};

\draw[ar] (a) -- (b)  node[lbl, midway, above] {\texttt{!rx\_empty}}
                      node[lbl, midway, below] {\texttt{rd=1}};
\draw[ar] (b) -- (op) node[lbl, midway, above] {\texttt{!rx\_empty}}
                      node[lbl, midway, below] {\texttt{rd=1}};
\draw[ar] (op) -- (s) node[lbl, midway, above] {\texttt{!rx\_empty}}
                      node[lbl, midway, below] {\texttt{rd=1}};
\draw[ar] (s.north) .. controls +(0,1.35) and +(0,1.35) .. (a.north)
  node[lbl, midway, above] {\texttt{!tx\_full} $\Rightarrow$ \texttt{wr=1} \quad (una operaci\'on completa)};

\foreach \n/\c in {a/{rx\_empty}, b/{rx\_empty}, op/{rx\_empty}, s/{tx\_full}} {
  \path[ar] (\n) edge[loop below, looseness=6, min distance=7mm]
    node[lbl, below=2pt] {\texttt{\c}} (\n);
}
\node[font=\sffamily\tiny, text=gris, anchor=north] at (5.55,-1.75)
  {los autoloops son la espera: aqu\'i la FSM pasa el 99,99 por ciento del tiempo};
\end{tikzpicture}
\end{center}
\vspace{4pt}

\begin{ficha}{Los tres estados de carga son id\'enticos}
Cada uno hace exactamente lo mismo: mirar \sg{rx\_empty}, y si hay algo, levantar
\sg{rd} y copiar \sg{r\_data} a su registro. Lo \'unico que cambia es
\emph{en qu\'e registro} guarda y \emph{a qu\'e estado} va. Los autoloops
(quedarse en el mismo estado mientras \sg{rx\_empty} est\'a en 1) son la espera:
la FSM se queda ah\'i, ciclo tras ciclo, sin hacer nada, durante los $\sim$52\,000
ciclos que tarda en llegar el byte siguiente.

\smallskip
De \sg{LOAD\_OP} solo se toman los 6 bits bajos: \sg{op\_next = r\_data[NB\_OP-1:0];}
Los dos bits altos del byte se descartan.
\end{ficha}

\subsection{Paso 7 --- La ALU: el bloque sin reloj}

Cuando el \sg{intf} pasa de \sg{LOAD\_OP} a \sg{SEND}, en ese mismo flanco se
cargan \sg{op\_reg} y ya estaban cargados \sg{a\_reg} y \sg{b\_reg}. Como la ALU
es combinacional, su salida se actualiza \textbf{sola}, sin ning\'un flanco de
clock, apenas cambian sus entradas.

\begin{intuicion}{La ALU como una calculadora siempre encendida}
No hay que ``disparar'' la ALU ni esperar que termine. Es un enredo de compuertas:
poneles A, B y OP en la entrada y unos nanosegundos despu\'es el resultado aparece
en la salida. Mientras A, B y OP no cambien, el resultado se queda ah\'i quieto
para siempre.

\smallskip
Por eso el \sg{intf} no necesita ning\'un estado de ``calcular''. Cuando llega a
\sg{SEND}, el resultado \emph{ya est\'a}: \sg{assign w\_data = i\_alu\_result;}
conecta la salida de la ALU directo a la entrada de la FIFO de transmisi\'on.
\end{intuicion}

\begin{ojo}{Lo \'unico que la ALU le exige al resto}
Que su retardo combinacional entre A/B/OP y el resultado entre en \textbf{un
per\'iodo de clock (10\,ns)}, porque ese resultado se captura en el flip-flop de
la FIFO en el flanco siguiente. Para una ALU de 8 bits en un Artix-7 esto sobra:
el camino cr\'itico es el sumador con acarreo, del orden de 2 a 3\,ns.

\smallskip
Si la ALU fuera de 64 bits o hiciera una multiplicaci\'on, este camino podr\'ia
no cerrar timing y habr\'ia que agregar un estado de espera o registrar la salida
en el medio (\emph{pipeline}).
\end{ojo}

% =========================================================================
\section{El viaje de vuelta: de la ALU a la PC}

\subsection{Paso 8 --- \texttt{SEND} y \texttt{fifo\_tx}}

En \sg{SEND}, si \sg{tx\_full} est\'a en 0, el \sg{intf} levanta \sg{wr} durante
un ciclo. En el flanco, la FIFO de transmisi\'on guarda el resultado y baja su
\sg{empty}. Aprovecha el mismo ciclo para enganchar el carry en un registro,
que es lo que despu\'es se ve en el LED 8.

\subsection{Paso 9 --- \texttt{uart\_tx} arranca solo}

Ac\'a hay un detalle lindo del dise\~no. El transmisor no recibe una orden de
nadie: mira si la cola tiene algo.

\begin{ficha}{El arranque autom\'atico}
\begin{center}
\sg{assign tx\_fifo\_not\_empty = \textasciitilde tx\_empty;}\\
\sg{uart\_tx u\_uart\_tx ( .tx\_start(tx\_fifo\_not\_empty), ... );}
\end{center}
Como \sg{uart\_tx} en \sg{IDLE} arranca apenas ve \sg{tx\_start==1}, poner ah\'i
directamente ``la cola no est\'a vac\'ia'' hace que el transmisor se dispare solo
y encadene bytes sin que nadie lo coordine. Y del otro lado:
\begin{center}\sg{.rd(tx\_done\_tick)}\end{center}
cuando el transmisor termina de sacar un byte, ese mismo pulso lo saca de la cola.
El lazo se cierra sin una sola l\'inea de control extra.
\end{ficha}

\subsection{Paso 10 --- El byte sale bit por bit}

El \sg{uart\_tx} recorre los mismos cuatro estados que el receptor, pero al
rev\'es: en vez de muestrear, \emph{maneja} la l\'inea. En \sg{START} la pone en 0
y espera 16 ticks; en \sg{DATA} saca \sg{b\_reg[0]} y despu\'es de 16 ticks
desplaza (\sg{b\_next = b\_reg >> 1}); en \sg{STOP} la pone en 1 durante 16 ticks
m\'as y levanta \sg{tx\_done\_tick}.

\begin{ojo}{El transmisor no se alinea con el tick al arrancar}
Cuando \sg{uart\_tx} pasa de \sg{IDLE} a \sg{START}, pone la l\'inea en 0
\textbf{inmediatamente}, en ese ciclo de clock --- y reci\'en despu\'es empieza a
contar los 16 ticks. Como el \sg{baud\_gen} corre libre, el primer tick puede
llegar entre 1 y 326 ciclos despu\'es. Resultado: el \textbf{bit de start puede
salir hasta 3,26 \textmu s m\'as corto} que los dem\'as (un 6 por ciento).

\smallskip
$\rightarrow$ ¿Rompe algo? No, porque el receptor del otro lado muestrea en el
\emph{centro} de cada bit y tiene medio bit de margen. Pero es exactamente el
tipo de detalle que vale la pena saber contestar si lo preguntan. Si se quisiera
evitar, bastar\'ia con condicionar la salida de \sg{IDLE} a \sg{if (tx\_start \&\&
s\_tick)}, a costa de hasta 3,26 \textmu s de latencia extra.
\end{ojo}

\begin{ficha}{El hueco entre tramas}
Cuando \sg{tx\_done\_tick} se levanta, en el mismo flanco pasan dos cosas: la FSM
vuelve a \sg{IDLE} y la FIFO saca el byte. Si qued\'o otro byte encolado,
\sg{tx\_empty} sigue en 0 y el transmisor arranca de nuevo al ciclo siguiente.
Es decir: entre dos tramas consecutivas hay \textbf{un ciclo de clock de hueco,
10\,ns}, contra 521\,600\,ns de trama. Despreciable, y adem\'as la l\'inea queda
en 1 (reposo), que es exactamente lo que espera el receptor.
\end{ficha}

% =========================================================================
\section{El cronograma completo}

Poner los tiempos juntos es lo que hace evidente qui\'en manda ac\'a.

\begin{ficha}{Una operaci\'on de punta a punta}
\renewcommand{\arraystretch}{1.25}
\begin{tabular}{@{}>{\sffamily}l r r p{5.9cm}@{}}
\toprule
\textbf{Etapa} & \textbf{Duraci\'on} & \textbf{En ciclos} & \textbf{Qui\'en trabaja} \\
\midrule
Trama de A            & 521,6 \textmu s  & 52\,160 & cable $+$ \sg{uart\_rx} \\
Trama de B            & 521,6 \textmu s  & 52\,160 & cable $+$ \sg{uart\_rx} \\
Trama de OP           & 521,6 \textmu s  & 52\,160 & cable $+$ \sg{uart\_rx} \\
\textbf{Trabajo del \sg{intf}} & \textbf{40 ns} & \textbf{4} & la FSM completa \\
C\'alculo de la ALU   & $\sim$3 ns       & $<1$    & combinacional \\
Trama del resultado   & 521,6 \textmu s  & 52\,160 & \sg{uart\_tx} $+$ cable \\
\midrule
\textbf{Total}        & \textbf{2,086 ms}& 208\,644& \\
\bottomrule
\end{tabular}
\end{ficha}

\begin{clave}
\textbf{El 99,998 por ciento del tiempo se lo lleva el cable.} La l\'ogica que
escribimos --- la FSM, la ALU, las FIFOs --- ocupa 4 ciclos de clock de los
208\,644 que dura la operaci\'on. Todo el dise\~no del UART existe para
adaptar un chip absurdamente r\'apido a un cable absurdamente lento, sin perder
un solo bit en el camino.
\end{clave}

\begin{intuicion}{D\'onde se va todo ese tiempo}
Si el sistema fuera una oficina: el cartero (el cable) tarda medio milisegundo
en traer cada sobre, y el empleado (la FPGA) resuelve el tr\'amite en 40
nanosegundos. El empleado pasa el 99,998 por ciento de su jornada mirando la
puerta. Los autoloops de las FSM --- \sg{rx\_empty} en \sg{LOAD\_A},
\sg{s\_tick==0} en \sg{DATA} --- son literalmente eso: mirar la puerta.
\end{intuicion}

% =========================================================================
\section{Tabla integradora: todas las se\~nales de un vistazo}

\begin{ficha}{Qui\'en genera qu\'e, y qui\'en la consume}
\renewcommand{\arraystretch}{1.2}
\begin{tabular}{@{}>{\ttfamily\small}l >{\sffamily\small}l >{\sffamily\small}l p{5.8cm}@{}}
\toprule
\textnormal{\textbf{\sffamily Se\~nal}} & \textbf{La genera} & \textbf{La usa} & \textbf{\sffamily Qu\'e significa} \\
\midrule
clk          & la placa    & todos         & el \'unico reloj del dise\~no \\
s\_tick      & baud\_gen   & rx, tx        & pulso de 1 ciclo, 16 por bit \\
rx\_sync     & el top      & uart\_rx      & la l\'inea serie ya sincronizada \\
rx\_done\_tick& uart\_rx   & fifo\_rx (wr) & ``termin\'e un byte'' \\
dout         & uart\_rx    & fifo\_rx      & el byte reci\'en armado \\
rx\_empty    & fifo\_rx    & intf          & ``no hay nada para leer'' \\
r\_data      & fifo\_rx    & intf          & cabeza de la cola (combinacional) \\
rd           & intf        & fifo\_rx      & ``sacame este byte'' \\
a, b, op     & intf        & ALU           & operandos registrados \\
result, carry& ALU         & fifo\_tx, LED & salida combinacional \\
wr           & intf        & fifo\_tx      & ``guardame el resultado'' \\
tx\_full     & fifo\_tx    & intf          & ``no entra nada m\'as'' \\
tx\_start    & fifo\_tx    & uart\_tx      & en realidad es \texttt{\textasciitilde tx\_empty} \\
tx\_done\_tick& uart\_tx   & fifo\_tx (rd) & ``sal\'i un byte entero'' \\
\bottomrule
\end{tabular}
\end{ficha}

% =========================================================================
\section{Las siete preguntas que m\'as se hacen}

\begin{enumerate}[leftmargin=1.6em, itemsep=7pt]

\item \textbf{¿Por qu\'e 16 ticks por bit y no 1?}\\
Con 1 tick por bit no hay forma de encontrar el centro: se muestrear\'ia en un
punto arbitrario, tal vez justo en la transici\'on. Con 16, el punto de muestreo
queda determinado con precisi\'on de $1/16$ de bit. Es el est\'andar de facto en
UARTs por hardware.

\item \textbf{¿Por qu\'e \sg{s\_reg} cuenta hasta 7 en \sg{START} y hasta 15 en
\sg{DATA}?}\\
7 (o sea 8 ticks) es \emph{medio} bit: lleva al centro del start. 15 (16 ticks) es
\emph{un} bit entero: mantiene el centrado en cada bit siguiente. Media m\'as
entera, entera, entera\dots

\item \textbf{¿El \sg{s\_tick} es un reloj?}\\
No. Es una \emph{habilitaci\'on}. Todos los flip-flops del dise\~no se disparan con
\sg{clk}; el tick aparece como una condici\'on \sg{if (s\_tick)} adentro de la
l\'ogica combinacional. Hay un \'unico dominio de reloj.

\item \textbf{¿Por qu\'e hace falta la FIFO si el \sg{intf} es 52\,000 veces m\'as
r\'apido que el cable?}\\
Porque ``m\'as r\'apido'' no es ``siempre disponible''. Si el \sg{intf} est\'a
trabado en \sg{SEND} justo cuando termina de llegar un byte, sin cola ese byte se
pierde. La FIFO desacopla \emph{disponibilidad}, no solo velocidad.

\item \textbf{¿Cu\'ando calcula la ALU?}\\
Siempre. No tiene reloj ni se\~nal de arranque. Su salida vale lo que valgan sus
entradas, permanentemente. El \sg{intf} no la ``dispara'': simplemente se asegura
de que A, B y OP est\'en cargados antes de leer el resultado.

\item \textbf{¿Qu\'e pasa si la PC manda 4 bytes en vez de 3?}\\
El cuarto queda en la FIFO y se toma como el \emph{A} de la operaci\'on siguiente.
El protocolo no tiene forma de resincronizarse: es posicional. Si se quisiera
robustez habr\'ia que agregar un byte de comando o un \emph{timeout} que vuelva a
\sg{LOAD\_A}.

\item \textbf{¿Por qu\'e todos los \sg{case} tienen \sg{default}?}\\
Por dos motivos distintos. En la l\'ogica combinacional, para que ninguna se\~nal
quede sin asignar en alg\'un camino --- eso sintetizar\'ia un \emph{latch}. Y en
las FSM, porque es la \emph{m\'aquina de estado segura} de la te\'orica: ante un
estado inv\'alido (un glitch, una part\'icula), el ciclo siguiente vuelve al
estado inicial en vez de quedar colgada.

\end{enumerate}

% =========================================================================
\section{Autoevaluaci\'on}

Si pod\'es contestar estas doce sin volver atr\'as, el flujo est\'a entendido.

\begin{enumerate}[leftmargin=1.6em, itemsep=4pt]
\item ¿Cu\'antos ciclos de clock dura un bit en el cable? ¿Y cu\'antos ticks?
\item ¿Qui\'en decide que empieza una trama, y c\'omo se entera?
\item Si cambiaras el baud rate a 115200, ¿qu\'e n\'umero tocar\'ias y cu\'anto
      valdr\'ia?
\item ¿Por qu\'e el shift register del receptor mete el bit nuevo por izquierda?
\item ¿Qu\'e pasar\'ia si en \sg{START} se contara hasta 15 en vez de hasta 7?
\item Un byte llega y \sg{rx\_empty} baja. ¿En qu\'e ciclo lo puede leer el
      \sg{intf}, y por qu\'e no antes?
\item ¿Por qu\'e \sg{r\_data} est\'a disponible sin esperar un ciclo despu\'es de
      \sg{rd}?
\item ¿Cu\'antos ciclos de clock tarda el \sg{intf} en procesar los tres bytes,
      una vez que ya llegaron todos?
\item ¿Por qu\'e la ALU no necesita una se\~nal de ``empezar'' ni una de ``listo''?
\item ¿Qui\'en le dice al \sg{uart\_tx} que tiene que transmitir?
\item ¿Cu\'anto tiempo de la operaci\'on completa es l\'ogica nuestra y cu\'anto es
      cable? Dalo en porcentaje.
\item ¿Por qu\'e la l\'inea \sg{i\_rx} pasa por dos flip-flops antes de entrar al
      receptor, y qu\'e le cuesta eso al dise\~no?
\end{enumerate}

\vfill
\begin{center}
\sffamily\footnotesize\color{gris}
TP2 --- Arquitectura de Computadoras --- FCEFyN, UNC
\end{center}

\end{document}
